\medskip

\textbf{Partie A}

\medskip



\[f(x) = 1 + x^2 - 2x^2\ln (x).\]



\medskip

\begin{enumerate}
\item ~
$\bullet~~$Limite en 0 : On sait que $\displaystyle\lim_{x \to 0} x^2\ln (x) = 0$, que $\displaystyle\lim_{x \to 0} x = 0$, donc par somme de limites :

$\displaystyle\lim_{x \to 0} f(x) = 1$.

$\bullet~~$Limite en $+ \infty$ : en écrivant $f(x) = 1 + x^2(1 - 2\ln (x))$, on a :

$\displaystyle\lim_{x \to + \infty}x^2 = + \infty, \:\displaystyle\lim_{x \to + \infty} - 2\ln (x) = - \infty$, puis $\displaystyle\lim_{x \to + \infty}1 - 2\ln (x) = - \infty$ et enfin par produit de limites :

$\displaystyle\lim_{x \to + \infty}f(x) = - \infty$.
\item 

Pour tout réel de l'intervalle $]0~;~+ \infty[, \: f'(x) = 2x - 4x\ln (x) - 2x^2 \times \dfrac1x = 2x - 4x\ln (x) - 2x = - 4x\ln (x)$.
\item 
Puisque $x \geqslant 0$, le signe de $f'(x)$ est l'opposé de celui de $\ln (x)$.

On sait que $\ln (x) < 0$ sur $[0~;~1[$, donc $f'(x) > 0$ sur $[0~;~1[$ et que 

$\ln (x) > 0$ sur $]1~;~+ \infty[$, donc $f'(x) < 0$ sur $]1~;~+ \infty[$.

La fonction $f$ est donc :

$\bullet~~$croissante sur [0~;~1] de 1 à $f(1) = 1 + 1 - 2 \times 1 \times 0 = 2$ ;

$\bullet~~$décroissante sur $[1~;~+\infty[$ de 2 à moins l'infini.

\begin{center}
\psset{unit=1cm,arrowsize=2pt 3}
\begin{pspicture}(7,2.5)
\psframe(7,2.5)\psline(0,1.5)(7,1.5)\psline(0,2)(7,2)\psline(1,0)(1,2.5)\psline(1.1,0)(1.1,2)
\uput[u](0.5,1.9){$x$} \uput[u](1.1,1.9){$0$} \uput[u](4,1.9){$1$} \uput[u](6.5,1.9){$+\infty$}
\uput[u](0.5,1.4){$f'(x)$} \uput[u](2,1.4){$+$} \uput[u](4,1.4){$0$} \uput[u](5.5,1.4){$-$}
\rput(0.5,0.75){$f(x)$}\uput[u](1.2,0){1}\uput[d](4,1.5){2} \uput[u](6.5,0){$- \infty$}
\psline{->}(1.5,0.4)(3.6,1.3)\psline{->}(4.4,1.3)(6.6,0.4)
\end{pspicture}
\end{center}

\item 
La fonction $f$ est continue car dérivable sur l'intervalle $[1~;~+ \infty]$ et décroissante de 2 à moins l'infini : d'après le théorème des valeurs intermédiaires, il existe un réel unique $\alpha$, avec $\alpha \in ]1~;~+ \infty[$ tel que $f(\alpha) = 0$.

Comme $f(\e) = 1 + \e^2(1 - 2\ln \e) = 1 + \e^2(1 - 2) = 1 - \e^2 \approx - 6,4$.

En appliquant le même théorème, on a donc $1 < \alpha < \e$.
\end{enumerate}

On admet dans la suite de l'exercice, que l'équation $f(x) = 0$ n'admet pas de solution sur l'intervalle ]0~;~1].

\emph{Rem. } : comme $0 \notin [1~;~2]$ le même théorème montre qu'il n'existe pas de réel $\beta \in [0~;~1]$ tel que $f(\beta) = 0$.

\begin{enumerate}[resume]
\item 










































On part de l'intervalle [1~;~2,7], (avec $\e \approx 2,7$)
{\ttfamily dichotomie(1)} donne par dichotomie un encadrement de $\alpha$ par deux réels $a$ et $b$ tels que $b - a \leqslant 10^{-1}$.

Comme $\alpha \approx 1,9$, C et D sont exclus et A ne donne pas un encadrement au dixième : reste la proposition B.
\end{enumerate}

\bigskip

\textbf{Partie B}

\medskip

On considère la fonction $g$ définie sur l'intervalle $]0~;~+ \infty[$, par 
\[g(x) = \dfrac{\ln (x)}{1 + x^2}.\]

On admet que $g$ est dérivable sur l'intervalle $]0~;~+ \infty[$ et on note $g'$ sa fonction dérivée.

On note $\mathcal{C}_g$ la courbe représentative de la fonction $g$ dans le plan rapporté à un repère \Oij.

\medskip

\begin{enumerate}
\item 
$g$ est un quotient de fonctions dérivables sur $]0~;~+ \infty[$, le dénominateur étant non nul (supérieur ou égal à 1), donc quel que soit $x \in ]0~;~+ \infty[$ :

$g'(x) = \dfrac{\dfrac 1x \times \left(1 + x^2\right) - 2x \ln (x)}{\left(1 + x^2\right)^2} = \dfrac{\dfrac{1 + x^2 - 2x^2\ln (x)}{x}}{\left(1 + x^2\right)^2}  = \dfrac{1 + x^2 - 2x^2\ln (x)}{x\left(1 + x^2\right)^2} = \dfrac{f(x)}{x\left(1 + x^2\right)^2}$.
\item 


Le résultat précédent montre que puisque $x > 0$ et $\left(1 + x^2\right)^2 > 0$, le signe de $g'(x)$ est celui du numérateur donc le signe de $f(x)$.

Or on a vu (Partie A question 3.) que $f(x) > 0$ sur $]0~;~\alpha[$ et $f(x) < 0$ sur $]\alpha~;~+ \infty[$.

La fonction $g$ est donc croissante sur $[0~;~\alpha]$, puis décroissante sur $[\alpha~;~+ \infty[$ avec un maximum $g(\alpha)$.

\emph{Rem.} : on sait que $f(\alpha) = 0 \iff 1 + \alpha^2 - 2\alpha^2 \ln \alpha \iff 2\alpha^2 \ln \alpha = 1 + \alpha^2 \iff $

$\ln \alpha = \dfrac{1 + \alpha^2}{2\alpha^2}$.

Donc $g(\alpha) = \dfrac{f(\alpha)}{1 + \alpha^2} = \dfrac{\dfrac{1 + \alpha^2}{2\alpha^2}}{1 + \alpha^2} = \dfrac{1}{2\alpha^2}$.
\item On note $T_1$ la tangente à $\mathcal{C}_g$ au point d'abscisse 1 et on note $T_{\alpha}$ la tangente à $\mathcal{C}_g$ au point d'abscisse $\alpha$.

$\bullet~~$On a $M(x~;~y) \in T_1 \iff y - g(1) = g'(1)(x - 1) \iff y - 0 = \dfrac12(x - 1) = \dfrac12 x - \dfrac12$ ;

$\bullet~~$On a vu que la fonction a un maximum en $\alpha$ et en ce point le nombre dérivé égal au coefficient directeur de la tangente en ce point (donc $T_{\alpha}$) est nul : la tangente est donc horizontale d'équation $y = g(\alpha)$ ou $y = \dfrac{1}{2\alpha^2}$. 


Donc $M(x~;~y) \in T_1 \cap T_{\alpha} \iff \dfrac12 x - \dfrac12 = \dfrac{1}{2\alpha^2}$, soit 

$x - 1 = \dfrac{1}{\alpha^2} \iff x = 1 + \dfrac{1}{\alpha^2}$ et $y = \dfrac{1}{2\alpha^2}$.

Donc $T_1 \cap T_{\alpha}\left(1 + \dfrac{1}{\alpha^2}~;~\dfrac{1}{2\alpha^2}\right)$.
\end{enumerate}

\bigskip
